\question[3]
Erkläre den Unterschied zwischen der \emph{Auflösung} und der \emph{Farbtiefe} 
eines Bildes und beschreibe jeweils, wie eine Verringerung dieser Größe
 zur Datenreduktion beiträgt.

\begin{solutionbox}{5cm}
Die Auflösung gibt an, aus wie vielen Bildpunkten (Pixeln) ein Bild besteht 
(z.\,B. Breite $\times$ Höhe). Wird die Auflösung verringert, werden mehrere benachbarte 
Pixel zu einem einzigen zusammengefasst -- es müssen weniger Pixelwerte gespeichert werden.
Die Farbtiefe gibt an, mit wie vielen Bits die Farbe bzw. 
der Helligkeitswert eines einzelnen Pixels codiert wird. Wird sie verringert, 
stehen pro Pixel weniger unterscheidbare Werte zur Verfügung, wodurch pro Pixel 
weniger Bits gespeichert werden müssen.
Beide Verfahren reduzieren die Datenmenge, führen jedoch zu einem Informationsverlust.
\end{solutionbox}